% Gerado por ferramentas/importar_relatorios.py a partir de MLP/ARTIGO_SIM_EMBARCADO-10.11.2026/ARTIGO_SIM_EMBARCADO.tex.
% Revisado à mão na dissertação (versão enxuta): edite este arquivo, não a fonte.
\subsection{Modelos e cadeia embarcada}
\label{sim:sec:sistema}

Os dados são as 56 corridas da Seção~\ref{sec:mm_bomba}; após a harmonização do limiar (canal A \(\ge 65{,}12\,\mathrm{dB_{AE}}\), canal B \(\ge 54{,}99\,\mathrm{dB_{AE}}\)), restam 281\,595 \textit{hits} em 3\,540 janelas de 5\,s. As duas redes têm os hiperparâmetros daquela seção: um regressor da frequência (36-64-32-1, 4\,481 parâmetros) e um classificador das válvulas (36-64-32-7, 4\,679 parâmetros). Para a implantação, ambas foram retreinadas com todas as janelas; o desempenho fora da amostra continua sendo o da Seção~\ref{sec:res_bomba}, e as acurácias ``no treino'' citadas adiante servem apenas para comparar implementações entre si.

Na cadeia embarcada (Figura~\ref{sim:fig:cadeia}), a lógica do MOMEA continua produzindo os registros. A MicroBlaze converte cada registro em unidades físicas, aplica o limiar harmonizado e acumula o \textit{hit} na janela corrente; a cada 5\,s, calcula os 36 descritores, padroniza e quantiza a entrada e executa as duas redes, em \textit{software} ou no acelerador HLS.

\begin{figure}[htbp]
	\centering
	\begin{tikzpicture}[
		font=\small, node distance=4mm and 4mm,
		bloco/.style={draw, rounded corners=2pt, align=center, minimum height=11mm, text width=22mm, fill=white},
		hw/.style={bloco, fill=blue!6},
		sw/.style={bloco, fill=orange!8},
		seta/.style={-{Latex[length=2mm]}, thick}]
		\node[hw] (det) {Detector de \textit{hits}\\(lógica do MOMEA)};
		\node[sw, right=of det] (conv) {Registro \(\to\)\\unidades físicas\\e limiar};
		\node[sw, right=of conv] (jan) {Janela de 5\,s:\\contagens, somas\\e percentis};
		\node[sw, right=of jan] (pre) {Imputação,\\padronização\\e int8};
		\node[hw, below=12mm of pre] (hls) {Núcleo int8\\(IP HLS)};
		\node[sw, left=of hls] (mbn) {Núcleo int8\\(C na MicroBlaze)};
		\node[bloco, left=of mbn, text width=30mm] (out) {Frequência (Hz)\\e classe de válvulas};
		\draw[seta] (det) -- node[above, font=\scriptsize]{32\,B} (conv);
		\draw[seta] (conv) -- (jan);
		\draw[seta] (jan) -- node[above, font=\scriptsize]{36} (pre);
		\draw[seta] (pre) -- (hls);
		\draw[seta] (pre.south west) to[out=-120, in=60] (mbn.north);
		\draw[seta] (mbn) -- (out);
		\draw[seta] (hls.south) to[out=-150, in=-30] (out.south);
	\end{tikzpicture}
	\caption{Cadeia embarcada. Em azul, lógica dedicada; em laranja, \textit{software} na MicroBlaze. O núcleo int8 tem uma única fonte em C, que compila para a MicroBlaze ou é sintetizada como IP pelo Vitis HLS.}
	\label{sim:fig:cadeia}
\end{figure}

\subsection{Implementação}
\label{sim:sec:implementacao}

\paragraph{Conversão e descritores.} O registro traz o pico em códigos (\(|x - z_c|\)), as contagens, o tempo de subida e a duração em amostras e a energia bruta \(\sum u^2\) em 48 bits. Com o fator de calibração do canal \(v_c\) (volts por código) e o ganho \(G = 10^{34/20}\),
\begin{equation}
	V_p = p\, v_c, \qquad
	A = 20 \log_{10}\frac{V_p}{G \cdot 1\,\mu\mathrm{V}}, \qquad
	E = \frac{\sum u^2\, (v_c/G)^2}{f_s \cdot 10^{-14}\,\mathrm{V^2 s}}, \qquad
	t = \frac{n}{f_s},
	\label{sim:eq:conversao}
\end{equation}
com \(E\) em eu e \(f_s = 10\,\)MSa/s; os fatores constantes são calculados uma vez por canal. Cada \textit{hit} atualiza contagens e somas da janela e, se não foi truncado, guarda sete valores para os percentis. No modo \textbf{exato}, esses valores ficam em vetores de até 1\,024 posições por canal, ordenados por \textit{heapsort} ao fechar a janela e interpolados como no \texttt{numpy} \cite{hyndman1996} (65,6\,kB); no modo \textbf{histograma}, cada variável tem um histograma de largura fixa (10,3\,kB e custo constante por \textit{hit}), alternativa para taxas altas de \textit{hits}.

\paragraph{Quantização.} A entrada é padronizada em ponto flutuante, saturada em \(\pm 8\) desvios e quantizada como \(q^x = \mathrm{rint}(z \cdot 127/8)\). Os pesos usam escala simétrica por camada (\(s_W = \max|W|/127\)), e as escalas das ativações foram calibradas no percentil 99,9 do máximo por janela. Cada camada densa com ReLU acumula em 32 bits e requantiza por uma multiplicação inteira e um deslocamento \cite{jacob2018,krishnamoorthi2018}:
\begin{equation}
	\begin{gathered}
		a_o = \sum_{i} q^{x}_{i}\, q^{W}_{io} + q^{b}_{o}, \qquad
		M = \mathrm{round}\!\left(2^{31}\,\frac{s_x\, s_W}{s_y}\right), \\
		q^{y}_{o} = \min\!\left(127,\ \left\lfloor \frac{\max(a_o, 0)\, M + 2^{30}}{2^{31}} \right\rfloor\right),
	\end{gathered}
	\label{sim:eq:requant}
\end{equation}
em que \(s_x\), \(s_W\) e \(s_y\) são as escalas da entrada, dos pesos e da saída. A camada de saída mantém o acumulador de 32 bits, convertido em Hz (regressor) ou usado no \(\arg\max\) (classificador). Como a Equação~\eqref{sim:eq:requant} é determinística, pode-se exigir igualdade bit a bit entre implementações; para isso, a quantização da entrada foi emulada em Python com as mesmas operações em precisão simples e o mesmo arredondamento do C \cite{goldberg1991,ieee754}, e o C foi compilado sem contração em FMA.

\paragraph{Núcleo int8 com fonte única.} As duas redes são escritas como laços de multiplicação e acumulação, com pesos em vetores constantes gerados do modelo treinado. O mesmo arquivo compila com \texttt{gcc} no computador, com \texttt{mb-gcc} para a MicroBlaze e no Vitis HLS como função de topo. No HLS, cada laço interno tem \textit{pipeline} com intervalo de iniciação 1, os pesos viram ROMs e a interface é AXI4-Lite; a escolha privilegia área, pois uma decisão a cada 5\,s não justifica paralelismo.

\subsection{Verificação por simulação}
\label{sim:sec:verificacao}

A verificação foi organizada em seis níveis, cada um com sua referência e seu critério de aceitação (Tabela~\ref{sim:tab:niveis}). Os níveis V1 a V3 usam todos os dados; V4 e V5 usam subconjuntos, porque a simulação RTL é lenta.

\begin{table}[htbp]
	\centering
	\small
	\caption{Níveis de verificação e critérios de aceitação.}
	\label{sim:tab:niveis}
	\begin{tabular}{L{0.06\textwidth}L{0.28\textwidth}L{0.30\textwidth}L{0.26\textwidth}}
		\toprule
		\textbf{Nível} & \textbf{O que se verifica} & \textbf{Como} & \textbf{Critério} \\
		\midrule
		V0 & Exportação do modelo & Rede em ponto flutuante reconstruída a partir dos pesos exportados, contra o \texttt{scikit-learn} & Diferença numérica desprezível; classes iguais \\
		V1 & Inferência em C no PC & 3\,540 janelas; float32 contra Python e int8 contra a emulação inteira & int8 bit a bit; float32 com diferença \(< 10^{-2}\)\,Hz \\
		V2 & Cadeia completa no PC & 281\,595 \textit{hits} \(\to\) descritores em C \(\to\) MLP int8, contra os descritores e as saídas int8 do Python & Erro normalizado \(\le 10^{-3}\); NaN idênticos; mesma classe e nível \\
		V3 & Núcleo HLS & csim e cosim C/RTL (Verilog, xsim) com 256 janelas & Acumuladores e classe idênticos \\
		V4 & Implementação do IP & Síntese e roteamento fora de contexto no XC7A35T, 100\,MHz & Temporização atendida \\
		V5 & Programa na MicroBlaze & Simulação RTL do processador executando o ELF final; até 2 janelas reais (163 \textit{hits}) e 2 vetores para a MLP & Resultados idênticos aos do PC; ciclos medidos \\
		\bottomrule
	\end{tabular}
\end{table}

O nível V5 substitui a medição em placa. Um sistema mínimo foi montado no Vivado 2022.2 para a mesma peça do MOMEA (xc7a35ticsg324-1L): MicroBlaze v11 com \textit{pipeline} de cinco estágios, sem caches, com \textit{barrel shifter}, multiplicador e divisor, 128\,kB de memória local em BRAM e um GPIO de 32 bits; foram simuladas as configurações sem FPU e com FPU estendida \cite{amd_ug984}. O programa escreve uma marca no GPIO antes e depois de cada etapa (processar os \textit{hits}, fechar a janela, executar a MLP), e o testbench registra o ciclo de cada mudança; ao fim de cada janela, escreve também a classe e parte do acumulador, e a sequência simulada tem de ser idêntica à do computador \cite{amd_ug900}. Como a memória é local e não há interrupções nem pilha de rede, os ciclos medidos são o melhor caso do código; no \textit{firmware} real, o tempo de parede será maior.
